% This is the general driver for the quick 
% NOVA-DAQ memo that details each of the dataformats
% that is used across the experiment
%
\documentclass[letterpaper]{article}
% \documentclass{article}
% \usepackage{hyperref}
%
% Use the ByteFields Package which makes nice 
% bytefield tables for our data diagrams
\usepackage{bytefield}
% Use the Vpage package which creates the proper 
% Page Layout (USletter, USlegal, etc...)
% see pages 88-89 of The LaTeX companion
%
\usepackage{vmargin}
%
% Use the doublespace package for linespacing
%\usepackage{doublespace}
% Use the setpace package for linespacing
\usepackage{setspace}
%
% Use both the AMS math and symbols packages 
% for equations
%
\usepackage{amsmath}
\usepackage{amssymb}
% Use the PS Tricks package for drawing
% postscript figures
%
% \usepackage{pstcol,pst-char}
\usepackage{pst-text}
\usepackage{pstricks}
%
% Use the color package
% \usepackage{xxcolor}
\usepackage{color}
\usepackage{array}
\usepackage{delarray}
%
% Use the standard Graphicx Package
% for graphics inclusion with the dvips driver
%
\usepackage[dvips]{graphicx}
\graphicspath{{.}{./figures/}}
%
% Use the EPS figure package for including 
% encapsulated postscript figures
%
%\usepackage{epsfig}
%\usepackage{psfig}
%
%
% Use the Xspace package for proper spacing 
% after macros
%
\usepackage{xspace}
%
% Use the PSFrag package for PS text replacement
%
\usepackage{psfrag}
%
% Use the Subfigure package to create arrays of 
% figures
%
\usepackage{subfigure}
%
% Use the Calc package for integer opps
%
\usepackage{calc}
%
% Use the TabularX package to create special tables
%
\usepackage{tabularx}
%
% use the Dcolumn package to allow for aligned columns
% in tables
%
\usepackage{dcolumn}
%
% Use the MakeIndex Package to generate an index
%
\usepackage{makeidx}
%%
% Use the enumerate package for custom lists
%
\usepackage{enumerate}
%
%
\usepackage{rotating}
% Set the Paper Size to Letter
%
\setpapersize{USletter}

%
%
\newcolumntype{L}{>{$}l<{$}}
\newenvironment{Cases}{\begin{array}\{{lL}.}{\end{array}}

%
% Set Up the margins to use 
%
\setmarginsrb{1in}{1in}{1in}{1in}{0pt}{10mm}{0pt}{5mm}
%
%
% $Id: general_macros.tex,v 1.1 2009/12/15 15:44:24 novacvs Exp $
%
% These are macros for use in the LaTeX documents
%
%
\newcommand{\Omit}[1]{}		% This command Omits the text placed in braces
%
% This Command produces the ampersand symbol, and should be used as a 
% replacement for the \& construct which is not handled well by search 
% and replace functions.
\newcommand{\amp}{\ensuremath{\&}}
%
% This command produces the percent sign, and should be used as a replacement
% for the \% construct which is not handled well by Emac's parser when color
% coding Text
\newcommand{\percent}{\ensuremath{\%}\xspace}
%
% This command produces the dollar sign, and should be used as a replacement
% for the \$ construct which is not handled well by Emac's parser when color
% coding Text
\newcommand{\dollar}{\ensuremath{\$}}

          % Include the general macros
% These are LaTeX Macro Definitions for common
% Kaon decays and symbols
%
% This macro file needs the Xspace command, so make
% sure to include the \usepackage{xspace} directive
% in your document
%
\newcommand{\Kzero}             {\ensuremath{{K}^{0}}\xspace}
\newcommand{\Kzerobar}          {\ensuremath{\bar{K}^{0}}\xspace}
\newcommand{\KzeroOne}          {\ensuremath{{K}_{1}^{0}}\xspace}
\newcommand{\KzeroTwo}          {\ensuremath{{K}_{2}^{0}}\xspace}
\newcommand{\Kshort}            {\ensuremath{{K}_{S}^{0}}\xspace}
\newcommand{\Klong}             {\ensuremath{{K}_{L}^{0}}\xspace}
\newcommand{\MuMu}              {\ensuremath{\mu^{+} \mu^{-}}\xspace}
\newcommand{\MuE}              {\ensuremath{\mu^{\pm} e^{\mp}}}
\newcommand{\EE}                {\ensuremath{e^{+} e^{-}}}
\newcommand{\PiPi}              {\ensuremath{\pi^{+} \pi^{-}}}
\newcommand{\ThreePi}           {\ensuremath{\pi^{+} \pi^{-} \pi^{0}}}
\newcommand{\MuMuEE}            {\ensuremath{\mu^{+} \mu^{-} e^{+} e^{-}}}
\newcommand{\FourLepton}        {\ensuremath{\ell^{+}\ell^{-}\ell^{+}\ell^{-}}}
\newcommand{\TwoLepton}         {\ensuremath{\ell^{+}\ell^{-}}}
\newcommand{\GammaGamma}        {\ensuremath{\gamma \gamma}}
\newcommand{\GammaGammaStar}    {\ensuremath{\gamma \gamma^{*}}}
\newcommand{\GammaStarGammaStar}{\ensuremath{\gamma^{*} \gamma^{*}}}
\newcommand{\KeThree}           {\ensuremath{\Klong \rightarrow \pi^{\pm} e^{\mp} \nu_{e}}\xspace}
\newcommand{\Ke}                {\ensuremath{K_{e3}}\xspace}
\newcommand{\Kmu}               {\ensuremath{K_{\mu 3}}\xspace}
\newcommand{\KmuThree}          {\ensuremath{\Klong\rightarrow\pi^{\pm} \mu^{\mp} \nu_{\mu}}\xspace}
\newcommand{\Kgamgam}           {\ensuremath{\Klong\rightarrow\GammaGamma}\xspace}
\newcommand{\Kgamgamstar}       {\ensuremath{\Klong\rightarrow\GammaGammaStar}\xspace}
\newcommand{\Kgamstargamstar}   {\ensuremath{\Klong\rightarrow\GammaStarGammaStar}\xspace}
\newcommand{\KtoMuMu}           {\ensuremath{\Klong\rightarrow\MuMu}\xspace}
\newcommand{\KtoMuE}            {\ensuremath{\Klong\rightarrow\MuE}\xspace}
\newcommand{\KtoMuMuEE}         {\ensuremath{\Klong\rightarrow\MuMuEE}\xspace}
\newcommand{\KtoEEEE}           {\ensuremath{\Klong\rightarrow\EE\EE}\xspace}
\newcommand{\KtoMuMuGamma}      {\ensuremath{\Klong\rightarrow\MuMu\gamma}\xspace}
\newcommand{\KtoEE}             {\ensuremath{\Klong\rightarrow\EE}\xspace}
\newcommand{\KtoEEGamma}        {\ensuremath{\Klong\rightarrow\EE \gamma}\xspace}
\newcommand{\KtoPiPi}           {\ensuremath{\Klong\rightarrow\PiPi}\xspace}
\newcommand{\KtoPiPiGamma}      {\ensuremath{\Klong\rightarrow\PiPi\gamma}\xspace}
\newcommand{\KtoPiPiEE}         {\ensuremath{\Klong\rightarrow\PiPi\EE}\xspace}
\newcommand{\KtoPiPiPi}         {\ensuremath{\Klong\rightarrow\ThreePi}\xspace}
%
\newcommand{\mutoe     }         {\ensuremath{\mu\rightarrow e}\xspace}
%
\newcommand{\bra}[1]            {\ensuremath{\left< #1 \right|}}
\newcommand{\ket}[1]            {\ensuremath{\left| #1 \right>}}
\newcommand{\CP}                {\ensuremath{\hat{C}\hat{P}}\xspace}
\newcommand{\Cabibbo}           {\ensuremath{\theta_{C}}\xspace}
\newcommand{\Cerenkov}          {\v Cerenkov\xspace}

\newcommand{\nova}              {NO$\nu$A\xspace}
%%% Local Variables: 
%%% mode: plain-tex
%%% TeX-master: "Thesis_Draft"
%%% End: 

%
\makeindex                      % Generate an index
%
\begin{document}                % Begin the document
%
% Reset the Counters
% \renewcommand{\thesection}{\arabic{section}}
% \renewcommand{\thefigure}{\arabic{section}.\arabic{subsection}.\arabic{figure}}
% \renewcommand{\theequation}{\arabic{section}.\arabic{subsection}.\arabic{equation}}
\newcommand{\bitlabel}[2]{%
  \bitbox[]{#1}{%
    \raisebox{0pt}[4ex][0pt]{%
      \turnbox{45}{\fontsize{7}{7}\selectfont#2}%
    }%
  }%
}%
%
% Start the Actual text
\title{\bf Determination of the Neutrino Time of Flight\\Fermilab to
  Ash River\\Nova-DocDB-11860}
% Put the Overall Section heading
\author{~A.Norman, FNAL}
\maketitle
%
%\tableofcontents               % Make a table of contents
\newpage
%
\normalsize                     % Return to Normal font size
\begin{spacing}{1.0}            % Double Space the Text
\section{Executive Summary - NO$\nu$A Timing Verification}

The timing of the NO$\nu$A detectors with respect to the NuMI beam spill was validated through multiple stages of calculations and experimental validation.  The general stages of the verification procedure proceeded as:

\begin{enumerate}
\item A priori calculation of accelerator signal propagation 
\item Measurement of the signal propagation to the Nova and Minos timing electronics
\item Observation of the neutrino rate excesses in the NOvA Surface detector
\item Comparison of the MINOS near detector timing to the NOvA near detector and accelerator timings
\item A priori calculation of the time of flight delays between the near and far NO$\nu$A detector locations 
\item Measurement by MINOS of the observed delays between the crossing of the beam at Minos near detector and the Minos far detector, as measured by the Minos master clock located at the Soudan Lab
\item Measurement by MINOS of the phase difference between the Minos master cesium clock and the GPS 1 second boundary
\item Measurement via a time transfer procedure using a rubidium clock between the MINOS cesium master clock and six independent GPS receiver systems located at Ash River
\item Measurement of the timing delays between the GPS based NOvA master clock and the clocks signals that are propagated to the detector for synchronization and readout 
\end{enumerate}

In particular the time transfer between the Soudan lab and the Ash River lab utilized a high precision cesium based atomic clock which was located at Soudan Lab.  The cesium atomic clock was used to train and phase align a portable rubidium based atomic clock.  The training was performed over a period of approximately 24 hours.  After the training of the rubidium clock, its phase was measured relative to the global positioning systems (GPS) one second boundaries signals at Soudan using a time interval counter with a precision for $\Delta T$ better than 100 ps.  The phase measurement was determined to have a Gaussian error with $\sigma=295ps$ as shown in Fig.~\ref{fig:minos_phase}.

%%%%%%%%%%%%%%%%%%%%%%%%%%%
% Begin Figure Block
\begin{figure}
\begin{center}
\includegraphics*[width=\textwidth*4/5]{rb_minos_alignment-2.eps}
\caption[Minos to Rubidium Phase]{Minos master cesium clock to trained rubidium clock phase offset after drift corrections.  Distribution is fit to a Gaussian profile for error determination.}
\label{fig:minos_phase}
\end{center}
\end{figure}
% End Figure Block
%%%%%%%%%%%%%%%%%%%%%%%%%%%

The rubidium clock was transported by automobile to the Ash River lab.  During this period the clock was powered by a large capacity uninterruptible power supply (battery backup) and monitored to ensure that clock was not interrupted during the trip.  The duration of the transport was roughly 2.5 hours.  Upon arrival at Ash River the rubidium clock was moved to the data acquisition computing facility where the NOvA timing systems are housed.  The phase difference between the GPS 1 second boundary provided by the NOvA timing system deployed at Ash River was then measured using the same time interval counter that was used at Soudan lab.  These phase intervals were measured for each of the three GPS receivers which had been in operations at the lab.  The time interval between the rubidium clock's reference pulses and the NOvA timing system's master clock was then measured utilizing the NOvA timing system's internal diagnostics.  

Three additional NOvA timing units were then deployed to Ash River lab and phases between these upgraded systems and the rubidium clock were measured in the same manners as the original three units.  Tests were then performed of the the power on and initialization characteristics of the timing system to determine if the timing phases were stable across power cycling, synchronization and re-initialization of the systems.

The rubidium clock was then transported back to Soudan Lab and the phase differences between the rubidium clock and the MINOS cesium master clock were remeasured to determine the total drift in the rubidium clock over the time during which the measurements were performed.  The clock was found to have a roughly linear drift of 0.45 ps/s accounting for a total accumulated drift of 16.5 ns 

\subsection{Summary}

The NOvA far detector timing system was determined to be correctly aligned to the GPS time standard.  The six independent timing distribution systems which were measured at Ash River were found to be aligned to the GPS clock to better than 16~ns per unit as shown in Fig.~\ref{fig:nova_time_align}.   Measurement errors and systematics related to the stability of the measurements across timing system restarts are included.

%%%%%%%%%%%%%%%%%%%%%%%%%%%
% Begin Figure Block
\begin{figure}
\begin{center}
\includegraphics*[scale=1]{nova_time_alignment.eps}
\caption[Nova Time Alignment]{Alignment of the NOvA timing distribution systems (TDUs) to the MINOS clock systems at Soudan Lab.  The measurements were corrected for drift of the rubidium atomic clock used for the time transfer and for known delays in the Soudan and Ash River timing systems.  The time deltas are reported as a function of the number of seconds from the start of the first measurement at Soudan Lab.}
\label{fig:nova_time_align}
\end{center}
\end{figure}
% End Figure Block
%%%%%%%%%%%%%%%%%%%%%%%%%%%

The result of the measurements were that the time of the crossing of the neutrino beam at the front face of the Ash River far detector was predicted to occur 2918.06 $\mu$s after the reception of the MIBS \$74 accelerator signal by the NOvA master timing system located at in the Minos service building and extrapolated to the far detector by the NOvA timing and synchronization system.

Neutrino candidates were observed in the far detector at a timing offset corresponding to 2918.5 $\mu$s after the reception of the NuMI trigger $t_0$ time.


\clearpage
\newpage
\section{Introduction}

The operation of the NOvA data acquistion (DAQ) and trigger hardware/software requires knowledge of the flight path of the neutrinos that are generated in the NuMI (Neutrinos at the Main Injector) and travel to the far detector.  Knowledge of this flight path allows for the computation of the propagation time required for the neutrinos to travel from the their origin to their crossing point with the NOvA far detector.  This propagation time is used to then set the global time delay that is used in the trigger system that determines the absolute start point of time window for which a ``beam'' trigger is issues.

The NuMI spill duration is nominally 10~$\mu$s and in baseline operations the trigger window would be set to be centered on this beam spill and include small side bands which allow for characterization of the detector in the interval surrounding the beam spill.  In the NOvA TDR these sidebands were proposed to be $\pm$10~$\mu$s which would yeild a total trigger window of 30~$\mu$s.  This trigger width is not optimal for the manner in which the DAQ systems were designed.  The design of the DAQ readout systems relys on a natural time based discretization of the data.  This discretization provides a smallest quanta of data which groups the hits into 50~$\mu$s long blocks.  As a result, the event builders are unable to break the data into windows shorter than 50~$\mu$s and requests for the extraction of data windows with a given width always result in the actual extraction of an integer number of contiguous 50~$\mu$s blocks.  As a result of this design the optimal setting for the trigger  window is to center the 10~$mu$s trigger inside of a 50~$\mu$s base window which then provides sidebands with a minimum width of 20~$\mu$s.  The trigger window then has the form of a half open interval:
\begin{equation}
\Delta T_{trig} = [ t_{0} + t^\prime_{delay} - 20~\mu s , t_{0} + t^{\prime}_{delay}+ 30~\mu s)
\end{equation}
Where $t_{0}$ is the time of the leading edge of the spill and $t^{\prime}_{delay}$ is the offset/delay between when the spill occured at FNAL and the time at which the beam crossed the detector.

When this trigger window is actually extracted from the data stream, it will be superimposed upon actual boundaries provided by the microslice quanta.  This will in all cases cause it [the trigger] to overlap with exactly 2 microslice blocks.  This is because the microslice intervals are considered closed:
\begin{equation}
\mu Slice = [t_n,t_n+50\mu s ]
\end{equation}
So that in the event that the front edge of a trigger window exactly coincides with one of the microslice boundaries, the trigger window is considered in overlap with both the trailing edge of first microslice and the leading edge of the second microslice.  This will cause both to be recorded into the data stream and for the final event on disk to have a width of 100~$\mu$s reguardless of the position of the trigger $t_0$ realative to the microslice boundary times.


\begin{figure}[h]
\begin{center}
\begin{bytefield}{10}
%  \bitheader[b]{}\\
  \bitlabel{1}{$t_0$}
  \bitbox{8}{\color{red}Trig $\Delta$t=50$\mu$s}
\end{bytefield}\\
\begin{bytefield}{20}
%  \bitheader[b]{}\\
  \bitbox{5}{\tiny $\mu$slice n}
  \bitbox{5}{\tiny $\mu$slice n+1}
  \bitbox{5}{\tiny $\mu$slice n+2}
  \bitbox{5}{\tiny $\mu$slice n+3}
%  \bitbox{5}{\color{lightgray}\rule{\width}{\height}}
 % \bitbox{5}{Pixel Adr}
 % \bitbox{7}{FEB Status}
 % \bitbox{6}{FEB ID}
 % \bitbox{1}{\tiny D \\ M}
 % \bitbox{1}{\tiny D \\ P}
 % \bitbox{1}{\tiny L \\ S}
\end{bytefield}
\caption[]{Event extracted from Microslice}
\end{center}
\end{figure}


In this scheme, the time $t^\prime_{delay}$ must be initially determined from an apriori calculation of the flight time of the neutrinos from the NuMI target to the NOvA detectors.  This document details how this calculation was performed and how certain inputs, such as the position of the NOvA detector were determines.


\section{Detector Positioning}

The determination of the position of the far detector was determined
using the information provided by the GPS receivers installed in the
timing system masters (TDU masters) at Ash River.  These three receivers
were located in the far detector computing center in relay rack 3 and
connected to a set of three antennas that are mounted on the roof of
the Ash River building.

The position of the near detector on the surface (NDOS) was determined
through the use of the master TDU mounted in the first electronics
rack at the NOvA surface building upper catwalk and connected to the
antenna mounted on the roof of the building.  Additionally the
position of the MINOS surface building was determined in a similar
manner using the TDUs mounted in the electronics racks in the MINOS
surface building and connected to two separate antennas mounted on a
mast on the roof of the building.

\subsection{GPS Antenna Position}

The Antennas on the Ash River building are mounted on a single mast
with four mount points for Antennas.  Each Antenna is mounted on a
mount point at a distance XXX from the central axis of the mast.  The
antennas are then connected to the recievers through an approximately
100 foot cable run that is common to all the antennas and receivers.

At the MINOS service building a similar antenna installations is used
with a single mast with two mount points and common cable runs down to
the GPS receivers.  For the NOvA surface building a single mast and
antenna were mounted and only a single cable runs down from the
antenna to the receiver.

\subsection{TDU NMEA Data Stream}

To extract the calculated GPS position, the receiver's NMEA data
stream output was connected to a serial console using a USB to serial
adapter (USB/serial chip on TDU).  The console was brought up using
the GNU ``screen'' utility using a connection speed of 38400 baud.

When this was started the MNEA data stream was continously sent to the serial
console in a fashion where it repleats its NMEA block 1 once every 5
seconds.  When extracted the blocks read similar to:

\begin{verbatim}
$GPGLL,4822.716417,N,09249.877638,W,210341.002,A,A*41
$GPRMC,210341.002,A,4822.716417,N,09249.877638,W,0.000,0.00,070313,,,A*40
$GPGGA,210341.002,4822.716417,N,09249.877638,W,1,12,0.71,382.823,M,-32.462,M,,*5D
$GPGSA,A,3,31,11,14,25,01,20,23,30,22,16,12,32,1.20,0.71,0.96*0C
$GPGSV,3,1,12,31,88,142,48,11,02,241,38,14,35,092,45,25,31,053,44*75
$GPGSV,3,2,12,01,20,256,44,20,40,300,45,23,11,285,42,30,27,171,44*74
$GPGSV,3,3,12,22,07,152,42,16,02,182,37,12,07,029,40,32,64,275,47*7C
$POLYT,210341.002,070313,421421.0022814,1730,421437.0022814,-281415,333.984,25,1127648768,00017,00017,*4C
$POLYP,210341.002,4822.716417,N,09249.877638,W,350.361,G3,001.9,005.1,0.000,0.00,0.000,,0.71,0.96,1.20,1.30,0.52,12,,0*23
$POLYS,12,31,U,142,88,48,255,11,U,241,02,38,255,14,U,092,35,45,255,25,U,053,31,44,255,01,U,256,20,44,255,20,U,300,40,45,255,23,U,285,11,42,255,30,U,171,27,44,255,22,U,152,07,42,255,16,U,182,02,37,201,12,U,029,07,40,255,32,U,275,64,47,255*73
$POLYI,JN,18,EXT,HPOS,VPOS,INT,TSYNC,*19 
\end{verbatim}

The \$GPGGA string is then used to determine the position of the fix.
In the case of the data above the \$GPGGA string is decode to the
position of the Ash River antenna array as shown in Table~\ref{table:gps_decode}.

\begin{table}[thb]
\begin{center}
\begin{tabular}{| c |r | r|}
\hline
Name                     & Data                      & Description \\
\hline
Sentence Identifier & \$GPGGA              & Global Positioning System Fix Data \\
Time                      &  210341.002         & 21:03:41.002 Z \\
Latitude                 & 4822.716417,N     & 48d 22.716417' N \\
Longitude              & 09249.877638,W  & 92d 49.877638' W \\
Fix Quality             & 1                           & 0=Invalid,
1=GPS fix, 2 = DGPS fix \\
Number of Satellites & 12                      & 12 satellites in view
and locked \\
Horizontal Dilution of Precision (HDOP) & 0.71 & final std error is
0.71 times input error \\
Altitude                   & 382.823,M           & 382.823 meters
above mean sea level \\
Height of geoid above WGS84 ellipsoid & -32.462,M & -32.462 meters \\
Checksum               & \*5D & Checksum for sentence \\
\hline
\end{tabular}
\caption[GPGGA Sentence Decoding]{\label{table:gps_decode} Decoding of the \$GPGGA sentence for
the Ash River GPS antenna array.}
\end{center}
\end{table}

The \$POLYP string can be similarly decoded to provide an alternative
position fix based on a user datum ellipsoid and the \$POLYT string
provides the GPS time information along with the errors in the
measurement and associated drifts.

\section{Flight Path Determination}

The flight path between the near detector and the far detector was
computing in two different manners using different assumptions for the
shape of the Eath.  These calculations are done both for self
consistancy and to determine the potential systemmatic error that is
introduced by the model of the Earth being used.

\subsection{Spherical Earth}

In the simplest approximation the Earth can be taken as a sphere with
a mean equitoral radius $R_E = 6,378,137$ meters.  In this representation
the longitude ($\theta$), latitude ($\phi$) and altitude (A) of a point can be converted to a
cartisan coordoinates in an Earth centered reference frame
representation.

The transformation $\{\theta, \phi, A\}$ to $\{p_x,p_y,p_z\}$ is given by:

\begin{align}
p_x &= ( R_E + A ) \sin(\theta)\cos(\phi) \\ 
p_y &= (R_E + A) \cos(\theta)\cos(\phi) \\
 p_z &= (R_E + A) \sin(\phi)
\end{align} 

In this case the longitude and latitude are taken from the GPS NMEA
data stream and the altitude that is used is the altitude above mean
sea level.

In the case of the near detector to far detector the parameters sets that were used
were:

\begin{table}
\begin{center}
\begin{tabular}{| c | c | c | c| c| }
\hline
Location & Latitude & Longitude & Altitude (m) & Offset Depth to Detector (m)\\
\hline
Ash River Ant-1 & 48 d 22.715780'' & 92 d 49.876930 & 348         & 12+7.8 \\  
Ash River Ant-2 & 48 d 22.715780'' & 92 d 49.876930 & 348.638  & 12+7.8 \\  
Ash River Ant-3 & 48 d 22.716416'' & 92 d 49.877639 & 350.346  & 12+7.8 \\  
\hline
Fnal Minos         &  41 d 50 m 23 s   & 88 d 16 m 13 s  & 228
& 105 + 15 \\
\hline
\end{tabular}
\caption[GPS Detector Locations]{\label{table:gps_det_loc} GPS
  coordinates of NOvA detectors used for flight path calculations.}
\end{center}
\end{table}

Because the values returned by the GPS receivers are the positions of
the antenna systems, offsets are applied to these reported positions
to account for the difference between the location of the antenna and
the location of the detector.  In each case the offset is computed to
correspond to the front face of the detector (face closest to the
beam) so that the flight time delays that are computed are for the
earliest arrival of particles between sites (i.e. delays are computing
from front face crossing to front face crossing and from that
corrected for position).  The depth offsets are computed to the
mid-point height of the detector (7.8~m for far det) and include the
approximate height of the detector over the top edge (12~m for far
det).  For the NuMI cavern the depth of the cavern is taken to be
105~m and the mount point of the antenna on the Minos service building
is taken as 15~m above the surface.

\begin{table}
\begin{center}
\begin{tabular}{| c | c | c | c| c| }
\hline
Location & $\Delta x_{offset}$ (m) & $Delta y_{offset}$ (m) & $Delta
z_{offset}$ (m)& Notes\\
\hline
Ash River Ant Base & TBD & TBD & TBD  & mount axis \\
Ash River Ant-1     & TBD & TBD & TBD  & south ant.\\  
Ash River Ant-2     & TBD & TBD & TBD  & north ant.\\  
Ash River Ant-3     & TBD & TBD & TBD  & west ant.\\  
\hline
Fnal NDOS         &  2 & 3   & -3 &  105 + 15 \\
\hline
\end{tabular}
\caption[Antenna to Detector
Offsets]{\label{table:antenna_det_offsets} Postions offsets used to
  correct for antenna verse detector front face positions.}
\end{center}
\end{table}

To comput these offsets the plane tangent to an earth centered sphere,
at the GPS measured antenna position is taken, and the vector between
antenna position and the center of the front face of the detector is
calculated.  The projects of the offset vector into the earth centered
fixed reference frame are extracted and the earth centered fixed
coordinates of the antenna are shifted accordingly.

The standard Euclidian metric is computed between the ECF coordinates
of the the near and far detectors to determine the straight line
flight path between the detectors.

\begin{align}
\Delta L &= \sqrt{\sum_{i} ( (p_i - \delta p_i) - (p_i^\prime - \Delta
  p_i^{prime})^2} \\
             &= 810,319~\text{m}
\end{align}

The neutrinos are assumes to have a speed equal to the speed of light,
c.  The resulting flight time is then:

\begin{align}
\Delta t &= \Delta L / c\\
             &= 2703.3~\mu\text{t}
\end{align}

\subsection{WGS84 Earth Model}

Due to the lattitudes and distances between the near and far detector
sites, a more accurate determination of the flight path can found by
using a deformed Earth model.  The standard description for the
deformation can be taken from the WGS84 model which establishes
elliptical major and minor axis for the earth and a set of
transformations from the ellipsoid into an Earth Centered Fixed (ECF)
representation.

The parameters used for the WGS84 model, as used by the both the NOvA
and Minos experiments are given by:

\begin{align}
R_a &= 6378137 ~\text{m}\\
R_b &= 6356752.31424518 ~\text{m}\\
\epsilon_a &= \sqrt{\frac{R^2_a - R^2_b}{R_a}}\\
\epsilon_b &= \sqrt{\frac{R^2_a - R^2_b}{R_a}}
\end{align}

With a prescription for transforming from the GPS reported ($\theta$,
$\phi$, A) coordinates into the ECF coordinates ($p_x$,$p_y$,$p_z$).

\begin{align}
R_{curv}(\theta) &= \frac{a}{\sqrt{1-\epsilon^2_a \sin \theta^2}} \\
p_x &= (R_{curv}(\theta) + A) \cos \theta \cos \phi \\
p_y &= (R_{curv}(\theta) + A) \cos \theta \sin \phi \\
p_z &= R^2_b/R^2_a (R_{curv}(\theta) + A) \sin \theta
\end{align}

In the case of the path length between the FNAL and Ash River timing
stations the path length is taken by comparing the distance between
the ECF coordinates of far detector and the Minos service building
position (offset to the beam depth and detector cavern).
\begin{align}
\Delta L &= \sqrt{\sum_{i} ( (p_i - \delta p_i) - (p_i^\prime - \Delta
  p_i^{prime})^2} \\
             &= 809,606~\text{m}
\end{align}

The time of flight is then computed assuming the propagation of the
beam at the speed of light

\begin{align}
\Delta t &= \Delta L / c \\
             &= 2700.56~\mu \text{s}
\end{align}

This time offsite is the delay between when the beam crosses the near
detector cavern and when the beam should cross the front face of the
far detector.
%\clearpage
%\section{Study - 4 Planes}

In the first study with the model, a geometry that only extended 4 planes deep was used.  This geometry was intended to verify the localized effect of staggering the planes and quantify the effects in terms of missed detector cells and total active path lengths.  The full module width was retained to minimize edge effects, and track trajectories were thrown at the model as describe previously.

The simulations were run with opening angle distributions extending to maximum angles of 15, 30 and 45 degrees to scan the available parameter space for detector crack effects in the forward direction, and other crack effects at larger angles dues to aliasing affect.

In a sample of 500,000 tracks with opening angles no greater than 5 degrees, 1309 tracks resulted in no detector cell registering a valid crossing in the unstaggered geometry.  In stagger geometries with 1/3, 1/2 and 2/3 cell width patterns, 0 tracks resulted in no detector cell registering a valid crossing in the staggered geometry.  These results were filtered to only include tracks that did not exist the detector volume prematurely.  The same analysis was performed for tracks with opening angles up to 10, 20 and 30 degrees and the results are tabulated in table~\ref{table:no-hits-4cell}.  Additionally, simulation was run with a geometry which allowed for shifting of detector modules by an ``assembly tolerance''.  In this simulation a maximum deviation of 4.1mm in the positioning of the first cell of a module was allowed, corresponding to a maximum accumulated deviation of 50cm in the position of the last module.  The deviations for these studies were generated uniformly, but only in the positive direction (i.e. no ``narrow'' modules were considered).  There geometries resulted in a significant reduction in the number of tracks that missed 4 successive cells.

%%%%%%%%%%%%%%%%%%%%%%%%
% Begin Table Block
\begin{table}[tb]
\begin{center}
\begin{tabular}{|l|c|c|c|c|c|}
\hline
Max Angle (deg)                 & No Stagger     & 1/3 Cell Stagger & 1/2 Cell Stagger & 2/3 Cell Stagger & Assembly\\
\hline
5                               & 1309           & 0                & 0                & 0                & 478  \\
10                              & 616            & 0                & 0                & 0                & 156  \\
20                              & 252            & 0                & 0                & 0                & 66   \\
30                              & 145            & 0                & 0                & 0                & 30   \\
\hline
\end{tabular}
\caption[4 Cell Missed Hits]{Number of track trajectories with no hits in a 4 plane deep cell geometry for different maximum track opening angles.  Base sample set was a collection of 500,000 track trajectories with random slopes generated uniformly from ($\theta,-\theta$) and with initial position uniformly generated across a full 384 cell detector.  Geometries considered are for different plane to plane staggerings, as well as the baseline ``aligned'' and ``aligned'' but with assembly tolerances.}
\label{table:no-hits-4cell}
\end{center}
\end{table}
% End Table Block
%%%%%%%%%%%%%%%%%%%%%%%%%%%%%%

Distributions of the total number of hit cells were then examined for both these geometries.  In each case both the number of hit cells on a given trajectory, and the total pathlength of active material for which that track passed through were measured.  These distributions are shown in figures \ref{fig:hitcells-4-10}-\ref{fig:pathlength-4-30}, where the unstaggered geometry is shown in blue and the staggered geometry's distribution in green on each plot.

%%%%%%%%%%%%%%%%%%%%%%%%%%%
% Begin Figure Block
\begin{figure}
\begin{center}
\includegraphics*[width=\textwidth*4/5]{hitcells-4cell-5-log.eps}
\caption[]{Number of hit cells for a 4 plane deep geometry at maximum incident angle of 5 degrees.  Unstaggered geometry is shown in blue and the half cell staggered geometry in green, and assembled geometry with tolerances is shown in black.  Magenta and Red distributions represent the staggering by 1/3 and 2/3 of a detector cell.}
\label{fig:hitcells-4-10}
\end{center}
\end{figure}
% End Figure Block
%%%%%%%%%%%%%%%%%%%%%%%%%%%
%%%%%%%%%%%%%%%%%%%%%%%%%%%
% Begin Figure Block
\begin{figure}
\begin{center}
\includegraphics*[width=\textwidth*4/5]{hitcells-4cell-30-log.eps}
\caption[]{Number of hit cells for a 4 plane deep geometry at maximum incident angle of 30 degrees.  Unstaggered geometry is shown in blue and the half cell staggered geometry in green.  The distribution in black represents the ``as assembled'' geometry, with magenta and red distributions corresponding to the 1/3 and 2/3 cell stagger patterns.}
\label{fig:hitcells-4-20}
\end{center}
\end{figure}
% End Figure Block
%%%%%%%%%%%%%%%%%%%%%%%%%%%

%%%%%%%%%%%%%%%%%%%%%%%%%%%
% Begin Figure Block
\begin{figure}
\begin{center}
\includegraphics*[width=\textwidth*4/5]{pathlength_4-5.eps}
\caption[]{Pathlength for track through active detector area for a 4 plane deep geometry at maximum incident angle of 5 degrees.  Unstaggered geometry is shown in blue, the half cell staggered geometry in green, and the ``as assembled'' in black.}
\label{fig:pathlength-4-10}
\end{center}
\end{figure}
% End Figure Block
%%%%%%%%%%%%%%%%%%%%%%%%%%%
%%%%%%%%%%%%%%%%%%%%%%%%%%%
% Begin Figure Block
\begin{figure}
\begin{center}
\includegraphics*[width=\textwidth*4/5]{pathlength_4-30.eps}
\caption[]{Pathlength for track through active detector area for a 4 plane deep geometry at maximum incident angle of 30 degrees.  Unstaggered geometry is shown in blue and the half cell staggered geometry in green, and the ``as assembled'' in black.}
\label{fig:pathlength-4-30}
\end{center}
\end{figure}
% End Figure Block
%%%%%%%%%%%%%%%%%%%%%%%%%%%

At small opening angles the staggering of detector planes results in an overall increase in both the number of cells that are traversed, and the total path length of active material that the tracks cross.  At angles between 10 to 20 degrees the staggering of the detector planes has the effect of on average slightly decreasing the number number of hit cells along a given trajectory due to a aliasing effect that results in a ``gap'' coming immediately after a long pathlength cell crossing.  This also has the effect of reducing the average pathlength of material that a given trajectory in this range traverses.  At maximum incident angles above 20 degrees there is essentially no net differences between the staggered and unstaggered geometries, except in the zero hit cells bin.

We can estimate that the effect of staggering the cells in terms of the elimination of 4 consecutive detectors cells, would be to regain 0.26\% of the very forward track angles (less than 5 degrees), but less than 0.03\% of all forward cone trajectories.  Staggering has almost no effect on the reduction of 3 of 4 or 2 of 4 type trajectories.

When the effect of random variation in the ``assembled'' position of the planes is considered the effect is less pronounced.  Staggering the planes by a 1/2 cell results in only a 0.09\% recovery of tracks in the very forward direction, and only a 0.006\% recovery of tracks over the forward cone.

%%% Local Variables: 
%%% mode: latex
%%% TeX-master: "nova-cellstagger"
%%% End: 

%\section{Study - 8 Planes}

In the second study with the model, a geometry that extended 8 planes deep in the Z direction was used.  This geometry was intended to examine the intermediate distance effect of staggering the planes and quantify the effects in terms of missed detector cells and total active path lengths, and tracks with significant numbers of missed detector cells.  The full module width was retained to minimize edge effects, and track trajectories were thrown at the model as describe previously.

The simulations were run with opening angle distributions extending to maximum angles of 5, 10, 20, 30, 45 degrees to scan the available parameter space for detector crack effects in the forward direction, and other crack effects at larger angles dues to aliasing affect as in the previous study.

In a sample of 500,000 tracks with opening angles no greater than 5 degrees, 580 tracks resulted in no detector cell registering a valid crossing in the unstaggered geometry, and 0 tracks resulted in no detector cell registering a valid crossing in all of the staggered geometry.  The ``as assembled'' geometry resulted in a significant reduction in the total number of trajectories that did not have hits.   These results are consistent with the actual ``cracks'' in the detector, and the incident trajectories of these tracks were mapped to show the sensitivity of the detector to these types of ``long''   These results are filtered to only include tracks that did not exit the detector volume prematurely.  The same analysis was performed for tracks with opening angles up to 10, 20 and 30 degrees and the results are tabulated in table~\ref{table:no-hits-8cell}.  In all cases, cell staggering of adjacent planes did eliminate trajectories which completely miss 8 sucessive cells.

%%%%%%%%%%%%%%%%%%%%%%%%
% Begin Table Block
\begin{table}[tb]
\begin{center}
\begin{tabular}{|l|c|c|c|c|c|}
\hline
Max Angle (deg)                 & No Stagger     & 1/3 Cell Stagger & 1/2 Cell Stagger & 2/3 Cell Stagger & Assembly\\
\hline
5                               &  580           & 0                & 0                & 0                & 7  \\
10                              &  310           & 0                & 0                & 0                & 9  \\
20                              &  139           & 0                & 0                & 0                & 3  \\
30                              &  58            & 0                & 0                & 0                & 7  \\
\hline
\end{tabular}
\caption[8 Cell Missed Hits]{Number of track trajectories with no hits in a 8 plane deep cell geometry for different maximum track opening angles.  Base sample set was a collection of 500,000 track trajectories with random slopes generated uniformly from ($\theta,-\theta$) and with initial position uniformly generated across a full 384 cell detector.  Geometries considered are for different plane to plane staggerings, as well as the baseline ``aligned'' and ``aligned'' but with assembly tolerances.}
\label{table:no-hits-8cell}
\end{center}
\end{table}
% End Table Block
%%%%%%%%%%%%%%%%%%%%%%%%%%%%%%

Distributions of the total number of hit cells were then examined for these geometries.  In each case both the number of hit cells on a given trajectory, and the total pathlength of active material for which that track passed through were measured.  These distributions are shown in figures \ref{fig:hitcells-8-10}-\ref{fig:pathlength-8-30}, where the unstaggered geometry is shown in blue and the staggered geometry's distribution in green on each plot.

%%%%%%%%%%%%%%%%%%%%%%%%%%%
% Begin Figure Block
\begin{figure}
\begin{center}
\includegraphics*[width=\textwidth*4/5]{hitcells-8cell-5-log.eps}
\caption[]{Number of hit cells for a 8 plane deep geometry at maximum incident angle of 5 degrees.  Unstaggered geometry is shown in blue and the half cell staggered geometry in green.  The distribution in black is the ``as assembled'' model and the magenta and red are the 1/3 and 2/3 cell staggererd geometries.}
\label{fig:hitcells-8-5}
\end{center}
\end{figure}
% End Figure Block
%%%%%%%%%%%%%%%%%%%%%%%%%%%
%%%%%%%%%%%%%%%%%%%%%%%%%%%
% Begin Figure Block
\begin{figure}
\begin{center}
\includegraphics*[width=\textwidth*4/5]{hitcells-8cell-30-log.eps}
\caption[]{Number of hit cells for a 8 plane deep geometry at maximum incident angle of 30 degrees.  Unstaggered geometry is shown in blue and the half cell staggered geometry in green.  The distribution in black is the ``as assembled'' model and the magenta and red are the 1/3 and 2/3 cell staggererd geometries.}
\label{fig:hitcells-8-10}
\end{center}
\end{figure}
% End Figure Block
%%%%%%%%%%%%%%%%%%%%%%%%%%%

One feature that is is important to notes, is that while the number of track trajectories with no registered hits decreases to zero in the staggered geometry, the number of tracks that are then produced with a number of hit cells less than the base number of planes increases slightly as a result of the staggered geometry's alignment of the center of a an active cell with an insensitive region in the next plane.  In the case of the very forward angles (less than 5 degrees) there is a roughly 8.4\% increase in the number of trajectories which result in these partial tracks in the staggered geometry, while recovering 0.1\% of missed trajectories.  This effect is shown in FIG.~\ref{fig:short_traj_8cell-5}.

%%%%%%%%%%%%%%%%%%%%%%%%%%%
% Begin Figure Block
\begin{figure}
\begin{center}
\includegraphics*[width=\textwidth*4/5]{shorttraj_comp-8cell-5.eps}
\caption[]{Angular distribution of track trajectories having a non-zero number of hit cells, but less cells than base number of detection planes (8).  Unstaggered geometry is shown in blue and the half cell staggered geometry in green.}
\label{fig:short_traj_8cell-5}
\end{center}
\end{figure}
% End Figure Block
%%%%%%%%%%%%%%%%%%%%%%%%%%%

The other effect of note is the total path length of material through which a track passes in each of the geometries.  These distributions are shown in FIG.~\ref{fig:pathlength-8-5}-\ref{fig:pathlength-8-30}.

%%%%%%%%%%%%%%%%%%%%%%%%%%%
% Begin Figure Block
\begin{figure}
\begin{center}
\includegraphics*[width=\textwidth*4/5]{pathlength_8-5.eps}
\caption[]{Pathlength for track through active detector area for a 8 plane deep geometry at maximum incident angle of 5 degrees.  Unstaggered geometry is shown in blue and the half cell staggered geometry in green.}
\label{fig:pathlength-8-5}
\end{center}
\end{figure}
% End Figure Block
%%%%%%%%%%%%%%%%%%%%%%%%%%%
%%%%%%%%%%%%%%%%%%%%%%%%%%%
% Begin Figure Block
\begin{figure}
\begin{center}
\includegraphics*[width=\textwidth*4/5]{pathlength_8-30.eps}
\caption[]{Pathlength for track through active detector area for a 8 plane deep geometry at maximum incident angle of 10 degrees.  Unstaggered geometry is shown in blue and the half cell staggered geometry in green.}
\label{fig:pathlength-8-5}
\end{center}
\end{figure}
% End Figure Block
%%%%%%%%%%%%%%%%%%%%%%%%%%%

Comparision of the distributions indicates that regular staggering of detector plans can eliminate the chance for a trajectory to pass through an exceedlingly long (8+ planes) path without registering in at least 1 cell.  It should be noted though that the same approximate suppression of these type of tracks can be obtained just by assuming the specified tolerances in the assembly of the detector planes.

%%% Local Variables: 
%%% mode: latex
%%% TeX-master: "nova-cellstagger"
%%% End: 

%input{study_100plane.tex}
%\clearpage
%\input{run.tex}
%
\end{spacing}
%
\newpage
% Put Bibliographic information last
%
\bibliography{nu_flight}    % Use the file ``thesis_biblio.bib'' for the biblio
\bibliographystyle{prsty}
%\bibliographystyle{apsrev}      % Indicate the sort by citation order style
%\bibliographystyle{apalike}     % This uses (name, year) as the citation style
%
%\begin{spacing}{1.0}
%\printindex                     % Print the Index
%\end{spacing}
%
\end{document}                  % End the document
